% Generated draft: requires mathematical and source review.
\documentclass[11pt]{article}
\usepackage[margin=25mm]{geometry}
\usepackage{fontspec}
\setmainfont{Latin Modern Roman}
\usepackage{amsmath,amssymb,mathtools,unicode-math}
\setmathfont{Latin Modern Math}
\usepackage{enumitem,needspace,xurl,hyperref}
\setlength{\parindent}{0pt}
\setlength{\parskip}{5pt}
\setlength{\emergencystretch}{3em}
\newcommand{\N}{\mathbb{N}}
\newcommand{\Z}{\mathbb{Z}}
\newcommand{\Q}{\mathbb{Q}}
\newcommand{\R}{\mathbb{R}}
\newcommand{\C}{\mathbb{C}}
\newcommand{\F}{\mathbb{F}}
\newcommand{\A}{\mathbb{A}}
\newcommand{\PP}{\mathbb{P}}
\newcommand{\E}{\mathbb{E}}
\newcommand{\eps}{\varepsilon}
\newcommand{\dd}{\,\mathrm{d}}
\begin{document}
\section*{UnsolvedMath ID: 30003460}
\textbf{Source problem code:} catalogue entry
\subsection{Definitions and mathematical statement}
Let \(X\) denote the intended ambient point space and let \(\mathcal L\) denote its collection of straight lines. The supplied record does not specify the ambient geometry or dimension; these remain unresolved parameters, rather than a claim about arbitrary incidence structures.

For an integer \(p\geq 3\), a finite family \(\mathcal F\subseteq\mathcal L\) has the \((p,3)\)-property if
\[
\forall\mathcal G\subseteq\mathcal F,\quad
|\mathcal G|=p\ \Longrightarrow\
\exists\ell_1,\ell_2,\ell_3\in\mathcal G\text{ pairwise distinct},\quad
\ell_1\cap\ell_2\cap\ell_3\ne\varnothing.
\]
Thus concurrency means that the three lines contain a common point of \(X\). A point set \(U\subseteq X\) pierces \(\mathcal F\) if
\[
\forall\ell\in\mathcal F,\qquad \ell\cap U\ne\varnothing.
\]
The supplied question asks whether, for every integer \(p\geq3\) and every such finite family with the \((p,3)\)-property, there exists a piercing set satisfying
\[
\exists U\subseteq X,\qquad |U|\leq p-2
\quad\text{and}\quad
\forall\ell\in\mathcal F,\ \ell\cap U\ne\varnothing.
\]
Here families are represented as sets of distinct lines; the source does not explicitly discuss repetitions. If \(|\mathcal F|<p\), the displayed \((p,3)\)-property is vacuous. The record supplies neither a restriction \(|\mathcal F|\geq p\) nor a separate convention for these small families, so the intended boundary formulation requires review. No such restriction has been inserted. The finite-family qualification comes from the supplied cleaned statement; the supplied original wording does not explicitly impose finiteness.
\subsection{Short English statement}
For each integer \(p\geq3\), must a finite family of straight lines in which every \(p\) lines contain three concurrent lines be pierced by at most \(p-2\) points? The ambient geometry and the treatment of families with fewer than \(p\) lines are unspecified and require clarification. This is a statement-only draft requiring human mathematical review.
\subsection{Sources}
\begin{itemize}
\item Local registry: \texttt{data/unsolvedmath/problems.json}, permanent UnsolvedMath ID \(30003460\), source problem code \texttt{catalogue entry}, title ``Piercing Line Families with the \((p,3)\)-Property.'' The supplied record belongs to the Oberwolfach Reports Open Problems collection. No individual proposer is supplied; the proposed year is recorded as 2017.
\item Supplied source citation: \emph{Discrete Geometry (2018)}, Oberwolfach Reports, pp.~1135--1205, DOI \texttt{10.4171/owr/2017/19}. Original Oberwolfach report: \url{https://doi.org/10.4171/owr/2017/19}. The background attributes the problem to mathematical workshop open-problem discussions documented in Oberwolfach Reports. The supplied 2018 citation and 2017 year metadata require reconciliation.
\item Supplied literature source: ``Ramsey type problems for lines and points,'' \url{https://www.renyi.hu/conferences/hunrus18/abstracts/grigoriev_abstract.pdf}. No author attribution is supplied in the record.
\item Literature checked at: \texttt{2026-08-21}. This is supplied dated curation, not an independently verified assessment. It labels the problem open, reports proofs in small cases such as \(p\leq6\), and describes the conjecture for all \(p\) as unresolved. The supplied background explicitly requests independent verification of time-sensitive status and cited claims.
\item The registry labels the statement \texttt{corrected\_verified} and says its mathematical content was preserved while notation, typography, and sentence structure were repaired. These are supplied editorial claims; neither the external sources nor the mathematical formulation have been independently verified here.
\end{itemize}
\end{document}
